\documentclass{article}
\usepackage[margin=0.5in]{geometry}
\usepackage{amsmath}
\usepackage{graphicx}
\usepackage{amssymb}
\usepackage{multicol}
\usepackage{xcolor}
\usepackage{amsthm}
\usepackage{mdframed}
\usepackage{tikz}
\usepackage{mathtools}
\usepackage [english]{babel}
\usepackage [autostyle, english = american]{csquotes}
\MakeOuterQuote{"}

\newcommand{\R}{\mathbb{R}}
\newcommand{\Z}{\mathbb{Z}}
\newcommand{\C}{\mathbb{C}}
\newcommand{\K}{\mathbb{K}}
\newcommand{\PP}{\mathbb{P}}
\newcommand{\E}{\mathcal{E}}
\newcommand{\e}{\vspace{3mm}}
\title{Práctica 2 - Cálculos combinatorios. Resultados igualmente probables.}
\author{Santiago}
\date{}
\begin{document}
    \maketitle
    \begin{enumerate}
        \item Consideremos mensajes enviados en código binario (0 y 1). ¿Cuántos mensajes es posible enviar con 10 dígitos o menos?\e\\
    Supongamos que puedo transmitir $n$ dígitos. Cada dígito en particular puede tener dos valores: 0 ó 1. Por ende, hay $2^n$ posibles mensajes posibles. Si consideramos que podemos enviar desde 1 hasta 10 dígitos, tenemos que \[\#\text{mensajes}=\sum\limits_{i=1}^{10}2^i=2^{11}-2=2046\]
        %\item Un experimentador está estudiando los efectos de la temperatura, la presión y el tipo de catalizador en la producción de cierta reacción química. Se consideran para las experiencias tres temperaturas diferentes, cuatro presiones distintas y cinco catalizadores diferentes. Si cualquier experimento particular implica utilizar una temperatura, una presión y un catalizador:
    \begin{enumerate}
        \item ¿Cuántos experimentos distintos son posibles realizar?\e\\
            Hay 3 temperaturas, 4 presiones y 5 catalizadores. Entonces,\[\#experimentos=3\cdot4\cdot5=60\]
        \item ¿Cuántos experimentos distintos existen que impliquen el uso de la temperatura más baja y las dos presiones más bajas?\e\\
            Hay una única temperatura más baja, y hay una restricción a 2 presiones en particular, por ende\[\#experimentos^\prime=1\cdot2\cdot5=10\]
        \item Suponga que se tiene que realizar cinco experimentos diferentes el primer día de experimentación y los experimentos se realizar al azar. ¿Cuál es la probabilidad de que se utilice un catalizador diferente en cada experimento?\e\\
            Los casos favorables: el primer experimento no tiene restricción, entonces hay 5 opciones. Para el segundo, no puedo usar el catalizador del primero, así que hay 4 opciones, por los mismos motivos, para el tercero hay 3, para el cuarto 2 y para el quinto 1. Casos totales: no hay restricción alguna en los experimentos, así que siendo $A=$"no se repite el catalizador", se tiene que\[P(A)=\frac{\text{casos favorables}}{\text{casos totales}}=\frac{5!}{5^5}\]
    \end{enumerate}
        %\item Se lanzan dos dados, sea "$a$" el número del primer dado y "$b$" el del segundo. ¿Cuál es la probabilidad de que la ecuación $x^2+ax+b^2=0$ tenga raíces reales?\e\\
    Para que las raíces sean reales:\[a^2-4b^2\geq0\to a\geq2b\]
    Tenemos que el espacio muestral es\[\Omega=\{(i,j):i,j=1,2,3,4,5,6\}\]
    Los casos favorables son \[B=\{(2,1)(3,1)(4,1)(5,1)(6,1)(4,2)(5,2)(6,2)(6,3)\}\]
    Sea $A=$"la ecuación tiene raíces reales"\[P(A)=\frac{\#B}{\#\Omega}=\frac{9}{36}\]
        %\item En una habitación, 10 personas tienen insignias numeradas del 1 al 10. Se eligen 3 personas al azar.
    \begin{enumerate}
        \item ¿Cuál es la probabilidad de que el número menor de las insignias de las personas elegidas sea 5?\e\\
            Los casos favorables involucran una selección de la insignia 5 y dos selecciones de números mayores. Hay $1\cdot5\cdot4$ casos favorables, mientras que los casos totales son $10\cdot9\cdot8$. Entonces\[P(\text{menor de las insignias es 5})=\frac{5\cdot4}{10\cdot9\cdot8}=\frac{20}{720}=\frac{1}{36}\]
        \item ¿Cuál es la probabilidad de que el número mayor de las insignias sea 5?\e\\
            Para los casos favorables necesito sacar el 5, luego tengo que sacar dos insignias menores. Con lo cual hay $1\cdot4\cdot3$ casos favorables.\[P(\text{mayor de las insignias es 5})=\frac{4\cdot3}{10\cdot9\cdot8}=\frac{12}{720}=\frac{1}{60}\]
    \end{enumerate}
        %\item Se eligen al azar dos números entre los primeros números $1,2,\dots,n$. ¿Cuál es la probabilidad que sean consecutivos si los escogemos con sustitución? ¿Y si lo hacemos sin sustitución?.\e\\
    Si es con sustitución, hay $n$ opciones para el primero, si se se seleccionó el 1 o $n$ hay 1 opción para el segundo, para los $n-2$ restantes hay 2. Los casos totales son $n^2$, entonces\[P(\text{consecutivos con sustitución})=\frac{(n-2)\cdot2+2\cdot1}{n^2}=\frac{2n-2}{n^2}=\frac{2}{n}-\frac{2}{n^2}\]
    El caso sin sustitución sólo me cambia los casos totales \[P(\text{consecutivos sin sustitución})=\frac{2n-2}{n\cdot(n-1)}=\frac{2}{n}\]
        %\item Cinco hombres y cinco mujeres son ordenados de acuerdo a su nota de examen. Asumimos que no hay dos notas iguales y que todos los posibles $10!$ reordenamientos son igualmente probables.
    \begin{enumerate}
        \item Encontrar la probabilidad de que una mujer obtenga el tercer puesto.\e\\
            En el tercer puesto puede haber tanto una mujer como un hombre, hay 5 casos favorables y 10 totales, entonces\[P(\text{mujer en el tercer puesto})=\frac{5}{10}=\frac{1}{2}\]
        \item Encontrar la probabilidad de que el puesto más alto alcanzado por una mujer sea el sexto.\e\\
            Para que esto sea posible, las mujeres deben ocupar los puestos 6,7,8,9 y 10. Por ende, tanto mujeres como hombres están restringidos a 5 posiciones.\[P(\text{mujeres a lo último})=\frac{5!5!}{10!}\]
    \end{enumerate}
        %\item En una canasta hay 30 manzanas de las cuales 5 están machucadas. Si elijo al azar, con orden y sin reemplazo 6 manzanas,
    \begin{enumerate}
        \item Defina un espacio muestral. ¿Cuántos elementos tiene?¿Es equiprobable?
        \item ¿Cuál es la probabilidad de que exactamente 3 de ellas estén machucadas?
        \item ¿Cuál es la probabilidad de que la tercera esté machucada?
    \end{enumerate}
        %\item En una fiesta se reparten al azar 10 caramelos entre 5 niños (un niño puede recibir más de un caramelo y obviamente otros niños pueden no recibir caramelos). Mi sobrino se encuentra en el grupo de niños. ¿Cuál es la probabilidad de que mi sobrino quede sin caramelo? (Ayuda: es conveniente suponer que tanto los niños como los caramelos están numerados).\e\\
    Se puede pensar como sigue\\
    %Insertar grafico
    $P(\text{no recibir})=1-P(recibir)=1-(P(\text{recibir 1})+\cdots+P(\text{recibir 10}))$
        %\item En una reunión hay 8 personas ¿Cuál es la probabilidad de que al menos 2 de ellas cumplan años el mismo día?¿Qué suposiciones son necesarias para el cálculo realizado?
    \[P(\text{al menos dos cumplan el mismo día})=1-P(\text{nadie comparta cumpleaños})\]
    Supongamos que ninguna persona nació en año bisiesto. Para los casos favorables, la primera persona no tiene restricción, así que hay 365 opciones, la segunda puede cumplir en 364, para evitar cumplir el mismo día que la primera, la tercera 363 y así hasta la octava. Para los casos totales, hay un total de $365^8$. Entonces\[P(\text{cumplir el mismo día})=1-\frac{365\cdot364\cdot363\cdot362\cdot361\cdot360\cdot359\cdot358}{365^8}=0.0743\]
        %\item Al tirar tres dados se puede obtener una suma de 9 de seis formas distintas, a saber: 126;135;144;225;234;333 y una suma de 10 también de 6 formas diferentes, a saber: 136;145;226;235;244;334. Sin embargo, la experiencia dice que es más fácil obtener 10 que 9 ¿Por qué?\\
Nota: este problema lo resolvió Galileo en el siglo $XVII$ a requerimiento de un jugador, el príncipe de Toscana.\e\\
    Sean los eventos $S_9=$"la suma es 9" y $S_{10}=$"la suma es 10". Tenemos que\[
    S_9=\begin{array}{l}
        \{126,162,216,261,612,621\\
        135,153,315,351,513,531\\
        144,414,441\\
        225,252,522\\
        234,243,324,342,423,432\\
        333\}
    \end{array}\qquad S_{10}=\begin{array}{l}
        \{136,163,316,361,613,631\\
        145,154,415,451,514,541\\
        226,262,622\\
        235,253,325,352,523,532\\
        244,424,442\\
        334,343,433\}
    \end{array}
    \]
    De donde se ve que $\#S_9=25$ mientras que $\#S_{10}=27$, con lo cual es más probable que salga el 10.
        %\item Los dados, tal y como los conocemos actualmente, se hicieron muy populares en la edad media. Chevalier De Mere propuso un enigma matemático: ¿Qué es más probable? ¿sacar al menos un 6 en cuatro tiradas con un sólo dado, o sacar al menos un doble seis en 24 tiradas con dos dados?\e\\
    Sea $A=$"sale al menos un 6" y $A_i=$"sale el 6 en el intento $i$". Entonces, \[P(A)=P(A_1\cup A_2\cup A_3\cup A_4)\]
    Por el principio de inclusión-exclusión\[P(A)=\sum\limits_{i=1}^4P(A_i)-\sum\limits_{i=1}^4\sum\limits_{j>i}^4P(A_iA_j)+\sum\limits_{i=1}^4\sum\limits_{j>i}^4\sum\limits_{k>j}^4P(A_iA_jA_k)-P(A_1A_2A_3A_4)\]
    En donde 
    \begin{align*}
        P(A_i)&=\frac{1\cdot5\cdot5\cdot5}{6^4}=\frac{5^3}{6^4}\quad\forall i\\
        P(A_iA_j)&=\frac{1\cdot1\cdot5\cdot5}{6^4}=\frac{5^2}{6^4}\quad\forall i\neq j\\
        P(A_iA_jA_k)&=\frac{1\cdot1\cdot1\cdot5}{6^4}=\frac{5}{6^4}\quad\forall i\neq j\neq k\\
        P(A_1A_2A_3A_4)&=\frac{1}{6^4}
    \end{align*}
    Entonces
    \[P(A)=\binom{4}{1}\frac{5^3}{6^4}-\binom{4}{2}\frac{5^2}{6^4}+\binom{4}{3}\frac{5}{6^4}-\binom{4}{4}\frac{1}{6^4}\]
    Ahora, para saber la probabilidad que de que haya al menos un doble 6 en 24 tiradas, se define el evento $B=$"sale al menos un doble 6", $P(B)=1-P(\text{no salga ningún doble 6})$.\[P(B)=1-\left(\frac{35}{36}\right)^{24}\]
    Faltaría comparar los resultados.
        %\item Tres matrimonios han comprado boletos para el teatro y se sientan en una fila formada por sólo seis asientos. Si toman sus asientos de un modo aleatorio:
    \begin{enumerate}
        \item ¿cuál es la probabilidad de que Juan y Paula (marido y mujer) se sienten en los dos asientos de la extrema izquierda?
        \item ¿cuál es la probabilidad de que Juan y Paula terminen sentados en forma contigua?
        \item ¿cuál es la probabilidad de que por lo menos una de las esposas quede sentada junto a su marido?
    \end{enumerate}
        %\item Se arrojan 5 dados distinguibles. Calcular la probabilidad de obtener:
    \begin{enumerate}
        \item escalera\e\\
            Hay dos maneras de hacer escalera: 1,2,3,4,5 ó 2,3,4,5,6.\[P(\text{escalera})=\frac{5\cdot4\cdot3\cdot2\cdot1\cdot2}{6^5}\]
        \item full
            \[P(\text{full})=\frac{6\x1\x1\x1\x1}{6^5}\]
        \item "pierna" ($aaabc$, siendo $a,b$ y $c$ distintos)
            \[P(\text{pierna})=\frac{6\x1\x1\x5\x4}{6^5}\]
        \item un par ($aabcd$)
            \[P(\text{par})=\frac{6\x1\x5\x4\x3}{6^5}\]
        \item par doble ($aabbc$)
            \[P(\text{doble par})=\frac{6\x1\x5\x1\x4}{6^5}\]
    \end{enumerate}
        %\item Dos amigos Luis y Pedro están con 6 personas y todos han colocado sus abrigos en un perchero. Al irse toman al azar un abrigo cada uno.
    \begin{enumerate}
        \item Hallar la probabilidad de que Luis tome su abrigo y Pedro no tome el suyo.\e\\
            Luis tiene una sóla opción para elegir su propio abrigo, mientras que Pedro no puede agarrar ni el suyo ni el de Luis, así que tiene 4 opciones. Una vez estos dos toman sus abrigos, los 4 restantes tienen libertad de elección, entonces \[P(\text{Luis toma su abrigo y Pedro no})=\frac{1\x4\x4!}{6!}\]
            Otra forma de verlo es con los eventos $A=$"Luis toma su abrigo", $B=$"Pedro toma su abrigo". Es de interés saber $P(A\cap B^c)$, entonces\[P(A\cap B^c)=P(A-A\cap B)=P(A)-P(A\cap B)=\frac{1\x5!}{6!}-\frac{1\x1\x4!}{6!}=\frac{4\x4!}{6!}\]
        \item ¿Cuál es la probabilidad que Luis o Pedro (al menos uno) escoja su propio abrigo?
            \[P(A\cup B)=P(A)+P(B)-P(A\cap B)=\frac{5!}{6!}+\frac{5!}{6!}-\frac{4!}{6!}=\frac{2\x5!-4!}{6!}\]
    \end{enumerate}
        %\item Una urna tiene 5 bolillas rojas, 6 azules y 8 verdes. Si se eligen 3 bolillas al azar. ¿Cuál es la probabilidad que:
    \begin{enumerate}
        \item todas sean del mismo color?
            \begin{center}
                $R=$"todas son rojas"\\
                $A=$"todas son azules"\\
                $V=$"todas son verdes"\\
                $B=$"todas del mismo color"
            \end{center}
            Como $R,A$ y $V$ son mutuamente excluyentes:\[P(B)=P(R)+P(A)+P(V)=\frac{5\x4\x3}{19\x18\x17}+\frac{6\x5\x4}{19\x18\x17}+\frac{8\x7\x6}{19\x18\x17}\]
        \item todas sean de distinto color?
            \[P(\text{todas distintas})=\frac{5\x6\x8}{19\x18\x17}\]
            ya que puedo agarrar entre 5 rojas, luego entre 6 azules y por último entre 8 verdes.
        \item Repetir los cálculos anteriores suponiendo que el muestreo se realiza con sustitución.\e\\
            Para que todas sean del mismo color ahora se tiene que \[P(B)=\frac{5^3}{19^3}+\frac{6^3}{19^3}+\frac{8^3}{19^3}\]
            Mientras que para que todas sean distintas\[P(\text{todas distintas})=\frac{5\x6\x8}{19^3}\]
    \end{enumerate}
        %\item En una ciudad hay 4 técnicos en reparación de televisores. Se han descompuesto 4 televisores.
    \begin{enumerate}
        \item Hallar la probabilidad que el técnico A sea llamado.\e\\
            Suponiendo que llaman al técnico de manera totalmente aleatoria, podemos plantear los siguientes eventos:\begin{center}
                $A_i=$"llaman al técnico A por el televisor $i$"\\
                $A=$"el técnico A es llamado al menos una vez"
            \end{center}
            Entonces\[P(A)=P(A_1\cup A_2\cup A_3\cup A_4)\]
            Por el principio de inclusión-exclusión:
            \[P(A)=\sum\limits_{i=1}^4P(A_i)-\sum\limits_{i=1}^4\sum\limits_{j>i}^4P(A_iA_j)+\sum\limits_{i=1}^4\sum\limits_{j>i}^4\sum\limits_{k>j}^4P(A_iA_jA_k)-P(A_1A_2A_3A_4)\]
            Como cada llamada es aleatoria \begin{align*}
                P(A_i)&=\frac{1}{4}\quad\forall i\\
                P(A_iA_j)&=\frac{1}{4^2}\quad\forall i,j\\
                P(A_iA_jA_k)&=\frac{1}{4^3}\quad\forall i,j,k\\
                P(A_1A_2A_3A_4)&=\frac{1}{4^4}
            \end{align*}
            Entonces,
            \[P(A)=\binom{4}{1}\frac{1}{4}-\binom{4}{2}\frac{1}{4^2}+\binom{4}{3}\frac{1}{4^3}-\binom{4}{4}\frac{1}{4^4}=\frac{175}{256}\]
        \item Hallar la probabilidad que exactamente $k$ técnicos sean llamados ($k=1,2,3,4$).\e\\
            Sea $B_k=$"$k$ técnicos son llamados" y $C_A=$"llaman únicamente al técnico A", se tiene que \[P(C_A)=\frac{1}{4^4}\]
            Esta probabilidad es la misma para el resto de técnicos, y como son eventos mutuamente excluyentes:\[P(B_1)=\binom{4}{1}\frac{1}{4^4}\]
            Supongamos ahora que quiero que llamen a dos en particular, sí o sí hay dos llamados que tienen que estar destinados a los involucrados, los otros dos pueden repartirse entre ambos de cualquier manera y hay que tener en cuenta que el orden de las llamadas entre estos dos puede ser cualquiera. Sea $D_{AB}=$"llaman al técnico A y B", entonces\[P(D_{AB})=\frac{1}{4}\x\frac{1}{4}\x\frac{2}{4}\x\frac{2}{4}\x\frac{4!}{2!2!}\]
            Y esta probabilidad es la misma para cualquier par que elija, por lo tanto\[P(B_2)=\binom{4}{2}\frac{4}{4^4}\]
            Para el caso de 3 se llega a\[P(B_3)=\binom{4}{3}\frac{3}{4^4}\]
            y para el de 4 \[P(B_4)=\binom{4}{4}\frac{1}{4^4}\]
    \end{enumerate}
        %\item En una fiesta entre 10 matrimonios (hombre-mujer), se eligen 4 de las mujeres y 4 de los hombres para formar un grupo. Hallar la probabilidad de que en ese grupo estén presentes ambos cónyuges de exactamente $k$ matrimonios, para $k=0,1,2,3,4$.\e\\
    Sean los eventos\[A_k=\text{"hay $k$ parejas"}\]
    Supongamos que arranco con $k=0$. Primero elijo 4 mujeres de entre las 10, $\binom{10}{4}$. Luego, tengo que seleccionar a 4 hombres de un grupo de 6, $\binom{6}{4}$, ya que hay 4 hombres que los tengo que descartar. Como el total de casos está dado por ${\binom{10}{4}}^2$, se tiene que\[P(A_0)=\frac{\binom{10}{4}\binom{6}{4}}{{\binom{10}{4}}^2}=\frac{1}{14}\]
    Para el caso de una pareja, puedo agarrar $\binom{10}{4}$ mujeres y puedo elegir entre $\binom{4}{1}$ para completar la pareja y luego $\binom{6}{3}$ para el resto de hombres\[P(A_1)=\frac{\binom{10}{4}\binom{4}{1}\binom{6}{3}}{{\binom{10}{4}}^2}=\frac{8}{21}\]
    El resto se argumenta de manera similar obteniendo\begin{align*}
        P(A_2)=\frac{\binom{10}{4}\binom{4}{2}\binom{6}{2}}{{\binom{10}{4}}^2}&=\frac{3}{7}\\
        P(A_3)=\frac{\binom{10}{4}\binom{4}{3}\binom{6}{1}}{{\binom{10}{4}}^2}&=\frac{4}{35}\\
        P(A_4)=\frac{\binom{10}{4}\binom{4}{4}\binom{6}{0}}{{\binom{10}{4}}^2}&=\frac{1}{210}
    \end{align*}
        %\item Se elije al azar un entero entre 1 y 100000. ¿Cuál es la probabilidad de que no sea múltiplo ni de 2 ni de 3 ni de 5?\e\\
    Se define $A=$"el número no es múltiplo de 2 ni 3 ni 5". Entonces $A^c=$"el número es múltiplo de 2 o 3 o 5". Sean $M_i=$"el número es múltiplo de i", se tiene que \[P(A)=1-P(A^c)=1-P(M_2\cup M_3\cup M_5)\]
    Por el principio de inclusión-exclusión\[P(M_2\cup M_3\cup M_5)=P(M_2)+P(M_3)+P(M_5)-P(M_6)-P(M_{10})-P(M_{15})+P(M_{30})\]
    De donde sigue que
    \[P(A^c)=\frac{1}{2}+\frac{33333}{100000}+\frac{1}{5}-\frac{16666}{100000}-\frac{1}{10}-\frac{6666}{100000}+\frac{3333}{100000}=0.73334\]
    Entonces\[P(A)=0.26666\]
        %\item Se forma una fila de 8 números al azar con dos números 1, dos números 2, dos números 3 y dos números 4. ¿Cuál es la probabilidad de que no haya números iguales consecutivos?\e\\
    Sea $A=$"no hay dos números iguales consecutivos", entonces $A^c=$"hay al menos un par consecutivo". Si defino $A_i=$"el número $i$ está consecutivo", se tiene que\[P(A^c)=P(A_1\cup A_2\cup A_3\cup A_4)\]
    Nuevamente se usará el principio de inclusión-exclusión\[P(A^c)=\sum\limits_{i=1}^4P(i)-\sum\limits_{i=1}^4\sum\limits_{j>i}^4P(A_iA_j)+\sum\limits_{i=1}^4\sum\limits_{j>i}^4\sum\limits_{k>j}P(A_iA_jA_k)-P(A_1A_2A_3A_4)\]
    Para la probabilidad de que el 1 salga consecutivo se puede pensar en el siguiente diagrama
    \begin{center}
        \begin{tikzpicture}
            \draw [thick, orange] (-2.75,0) -- (-3.75,0);
            \draw [thick, orange] (-1.5,0)  -- (-2.5,0);
            \draw [thick, orange] (-0.25,0)  -- (-1.25,0);
            \draw [thick, orange] (0,0)  -- (1,0);
            \draw [thick, orange] (1.25,0)  -- (2.25,0);
            \draw [thick, orange] (2.5,0)  -- (3.5,0);
            \draw [thick, orange] (3.75,0)  -- (4.75,0);
            \filldraw[color=red!60, fill=red!5, very thick](-3.25,0.25) circle (0.2);
            \node at (-3.25,0.25) {$1$};
        \end{tikzpicture}
    \end{center}
    En donde el 1 simboliza el par 1 y los 6 lugares restantes son para el resto de los números, entonces\[P(A_1)=\frac{\frac{7!}{2!2!2!}}{\frac{8!}{{2!}^4}}=\frac{1}{4}\]
    y esto es igual para cada $A_i$. Para el caso de que dos pares sean consecutivos, se procede de manera similar:
    \begin{center}
        \begin{tikzpicture}
            \draw [thick, orange] (-2.75,0) -- (-3.75,0);
            \draw [thick, orange] (-1.5,0)  -- (-2.5,0);
            \draw [thick, orange] (-0.25,0)  -- (-1.25,0);
            \draw [thick, orange] (0,0)  -- (1,0);
            \draw [thick, orange] (1.25,0)  -- (2.25,0);
            \draw [thick, orange] (2.5,0)  -- (3.5,0);
            \filldraw[color=red!60, fill=red!5, very thick](-3.25,0.25) circle (0.2);
            \node at (-3.25,0.25) {$1$};
            \filldraw[color=red!60, fill=red!5, very thick](-2,0.25) circle (0.2);
            \node at (-2,0.25) {$2$};
        \end{tikzpicture}
    \end{center}
    \[P(A_1A_2)=\frac{\frac{6!}{2!2!}}{\frac{8!}{2!^4}}=\frac{1}{14}\]
    Esto también se cumple para cualquier par. En el caso de tres, se llega a\[P(A_iA_jA_k)=\frac{\frac{5!}{2!}}{\frac{8!}{2!^4}}=\frac{1}{42}\]
    Y Finalmente para que los cuatro sean consecutivos\[P(A_1A_2A_3A_4)=\frac{4!}{\frac{8!}{2!^4}}=\frac{1}{105}\]
    Por lo tanto\[P(A)=1-\left[\binom{4}{1}\frac{1}{4}-\binom{4}{2}\frac{1}{14}+\binom{4}{3}\frac{1}{42}-\binom{4}{4}\frac{1}{105}\right]=\frac{12}{35}\]
    \end{enumerate}
\end{document}